%------------------------------------------------------------------------------%
\documentclass[a4paper,12pt,fleqn]{article}

\usepackage{calculo_varias_variaveis-1}

\ShowAnswers

%------------------------------------------------------------------------------%
\begin{document}

\makeHeader{CFVVI}{Prova 4 -- Turma 7}{24/11/2025}

% 01
%------------------------------------------------------------------------------%
\exercise{50\%}
Determine e classifique os pontos críticos da função
\(
  f(x,y) = e^{xy} - x^2 - y^2
\)

\clearpagequestiononly

\begin{answer}
  Calculando o gradiente
  \[
    f_x(x, y)
    = y e^{xy} - 2x
  \]
  \[
    f_y(x, y)
    = x e^{xy} - 2y
  \]
  Igualando o gradiente ao vetor nulo
  \[
    \begin{cases}
      y e^{xy} - 2x = 0\\
      x e^{xy} - 2y = 0
    \end{cases}
  \]
  rearranjando
  \[
    \begin{cases}
      2x = y e^{xy} \\
      2y = x e^{xy}
    \end{cases}
  \]
  Subtraindo as equações
  \[
    2(x-y) = -(x-y)e^{xy}
  \]

  Caso 1: $x\neq y$, dividimos os dois lados por $x-y$, obtendo
  \(
    2 = -e^{xy}
  \)
  que não possui solução.

  Caso 2: $x=y$, a equação \(2(x-y) = -(x-y)e^{xy}\) é satisfeita

  Substituindo na primeira equação e obtemos
  \begin{align*}
    2x & = y e^{xy} \\
    2x & = x e^{x^2}
  \end{align*}

  Se $x=0$, essa equação é satisfeita e temos o primeiro ponto crítico \(P_1=(0, 0)\)

  Se $x\neq y$, dividimos os dois termos por $x$
  \begin{align*}
    2x      & = x e^{x^2}        \\
    2       & = e^{x^2}          \\
    \ln(2)  & = x^2              \\
    \abs{x} & = \sqrt{\ln(2)}    \\
    x       & = \pm\sqrt{\ln(2)}
  \end{align*}
  Encontramos os pontos críticos
  \(P_2 = \left( \sqrt{\ln(2)},  \sqrt{\ln(2)}\right)\) e
  \(P_3 = \left(-\sqrt{\ln(2)}, -\sqrt{\ln(2)}\right)\)

  \bigskip
  Calculando a Hessiana de $f$
  \[
    f_{xx}(x, y)
    = y^2 e^{xy}-2
  \]
  \[
    f_{yy}(x, y)
    = x^2 e^{xy}-2
  \]
  \[
    f_{xy}(x, y)
    = xy e^{xy}+e^{xy}
  \]
  Discriminante (Determinante da Hessiana)
  \[
    D(x, y) = f_{xx}(x, y)f_{yy}(x, y) - f^2_{xy}(x, y)
  \]

  \bigskip
  Classificando o ponto \((0, 0)\)
  \[
    f_{xx}(0, 0) = -2
  \]
  \[
    f_{yy}(0, 0) = -2
  \]
  \[
    f_{xy}(0, 0) =  1
  \]
  \[
    D(0, 0) = 4-1 = 3 > 0
  \]
  Como \(f_{xx}(0, 0) = -2 < 0\), \textbf{o ponto \((0,0)\) é um máximo local}

  \bigskip
  No ponto \(\left(\sqrt{\ln 2}, \sqrt{\ln 2}\right)\)
  \[
    f_{xx}\left(\sqrt{\ln 2}, \sqrt{\ln 2}\right)
    = 2\ln 2 -2
  \]
  \[
    f_{yy}\left(\sqrt{\ln 2}, \sqrt{\ln 2}\right)
    = 2\ln 2 -2
  \]
  \[
    f_{xy}\left(\sqrt{\ln 2}, \sqrt{\ln 2}\right)
    = 2\ln 2 +2
  \]
  \[
    D\left(\sqrt{\ln 2}, \sqrt{\ln 2}\right)
    = (2\ln 2 -2)^2 - (2\ln 2 +2)^2
    < 0
  \]
  Portanto, é um \textbf{ponto de sela}

  \bigskip
  No ponto \(\left(-\sqrt{\ln 2}, -\sqrt{\ln 2}\right)\)
  \[
    f_{xx}\left(-\sqrt{\ln 2}, -\sqrt{\ln 2}\right)
    = 2\ln 2 -2
  \]
  \[
    f_{yy}\left(-\sqrt{\ln 2}, -\sqrt{\ln 2}\right)
    = 2\ln 2 -2
  \]
  \[
    f_{xy}\left(-\sqrt{\ln 2}, -\sqrt{\ln 2}\right)
    = 2\ln 2 +2
  \]
  \[
    D\left(-\sqrt{\ln 2}, -\sqrt{\ln 2}\right)
    = (2\ln 2 -2)^2 - (2\ln 2 +2)^2
    < 0
  \]
  Portanto, é um \textbf{ponto de sela}

\end{answer}

% 02
%------------------------------------------------------------------------------%
\exercise{50\%}
Determine os valores máximo e mínimo absolutos da função
\(
  f(x, y) = y\sin(x)
\)
na região
\(
  R = \{(x,y)\in\R^2\colon -1 \le x \le 2,\, -2 \le y \le 2\}
\)

\textit{Dica: Considere as aproximações
\(\sin(1)\approx 0,8\),
\(\sin(2)\approx 0,9\) e
\(\sin(3)\approx 0,1\)}

\clearpagequestiononly

\begin{answer}
  \textbf{Interior}

  Derivadas parciais
  \[
    f_x = y\cos(x) \qquad\qquad
    f_y =  \sin(x)
  \]
  Pontos críticos no interior
  \[
    y\cos(x) = 0 \qquad\qquad
     \sin(x) = 0
  \]
  Impondo \(\sin(x)=0\) em \(x\in[-1, 2]\) temos $x=0$.
  A primeira equação se torna
  \begin{align*}
    y\cos(x) & = 0 \\
    y\cos(0) & = 0 \\
    y        & = 0
  \end{align*}
  O ponto critico no interior da região é \((0, 0)\), onde função vale
  \[
    f(0, 0)
    = 0\sin(0)
    = 0
  \]

  \textbf{Fronteira 1}: onde \(x=-1\) e \(-2\le y \le 2\)
  \[
    g(y)
    = f(-1, y)
    = y \sin(-1)
    = -y \sin(1)
  \]
  Como \(g'(y) = -\sin(1) \neq 0\)
  a função $g(y)$ não possui pontos críticos,
  consideramos apenas
  os extremos do intervalo $y=-2$ e $y=2$
  \begin{gather*}
    f(-1, -2)
    = g(-2)
    = 2\sin(1)
    \\
    f(-1, 2)
    = g(2)
    = -2\sin(1)
  \end{gather*}

  \textbf{Fronteira 2}: onde \(x=2\) e \(-2\le y \le 2\)
  \[
    h(y)
    = f(2, y)
    = y\sin(2)
  \]
  Como \(h'(y) = \sin(2) \neq 0\)
  a função $h(y)$ não possui pontos críticos,
  consideramos apenas
  os extremos do intervalo $y=-2$ e $y=2$
  \begin{gather*}
    f(2, -2)
    = h(-2)
    = -2\sin(2)
    \\
    f(2, 2)
    = h(2)
    = 2\sin(2)
  \end{gather*}

  \textbf{Fronteira 3}: onde \(y=-2\) e  \(-1 \le x \le 2\)
  \[
    p(x)
    = f(x, -2)
    = -2\sin(x)
  \]
  Como \(p'(x) = -2\cos(x)\), impondo \(p'(x)=0\) com \(-1 \le x \le 2\)
  temos que \(x=\frac{\pi}{2}\) é um ponto crítico no interior do intervalo.
  Consideramos os pontos \(x=-1\), \(x=\frac{\pi}{2}\) e \(x=2\)
  \begin{gather*}
    f(-1, -2)
    = p(-1)
    = -2\sin(-1)
    = 2\sin(1)
    \\
    f\left(\sfrac{\pi}{2}, -2\right)
    = p\left(\sfrac{\pi}{2}\right)
    = -2 \sin\left(\sfrac{\pi}{2}\right)
    = -2
    \\
    f(2, -2)
    = p(2)
    = -2\sin(2)
  \end{gather*}

  \textbf{Fronteira 4}: onde \(y=2\) e  \(-1 \le x \le 2\)
  \[
    q(x)
    = f(x, 2)
    = 2\sin(x)
  \]
  Como \(q'(x) = 2\cos(x)\), impondo \(q'(x)=0\) com \(-1 \le x \le 2\)
  temos que \(x=\frac{\pi}{2}\) é um ponto crítico no interior do intervalo.
  Consideramos os pontos \(x=-1\), \(x=\frac{\pi}{2}\) e \(x=2\)
  \begin{gather*}
    f(-1, 2)
    = q(-1)
    = 2\sin(-1)
    = -2\sin(1)
    \\
    f\left(\sfrac{\pi}{2}, 2\right)
    = q\left(\sfrac{\pi}{2}\right)
    = 2 \sin\left(\sfrac{\pi}{2}\right)
    = 2
    \\
    f(2, 2)
    = q(2)
    = 2\sin(2)
  \end{gather*}

  Coletando todos os candidatos a extremo absoluto temos
  \[
    \begin{array}{rrrr}
      \hline
       x             &  y  &  f(x, y)  & \approx      \\ \hline
       0             &  0  &           &   0,0 \\
      -1             & -2  &  2\sin(1) &   1,6 \\
      -1             &  2  & -1\sin(1) &  -0,8 \\
       2             & -2  & -2\sin(2) &  -1,8 \\
       2             &  2  &  2\sin(2) &   1,8 \\
      -1             & -2  &  2\sin(1) &   1,6 \\
      \sfrac{\pi}{2} & -2  &           &  -2,0 \\
       2             & -2  & -2\sin(2) &  -1,8 \\
      -1             &  2  & -2\sin(1) &   1,6 \\
      \sfrac{\pi}{2} &  2  &           &   2,0 \\
       2             &  2  &  2\sin(2) &  -1,8 \\ \hline
    \end{array}
  \]

  O mínimo absoluto é $-2$ em \(\left(\frac{\pi}{2}, -2\right)\)

  O máximo absoluto é 2 em \(\left(\frac{\pi}{2}, 2\right)\)
\end{answer}

%------------------------------------------------------------------------------%
\end{document}
%------------------------------------------------------------------------------%


